\chapter{Drowning}

\section*{Definition and Overview}
Drowning is the process of experiencing respiratory impairment from being submerged or immersed in liquid, and it can be fatal or non-fatal. Drowning is a leading cause of accidental death worldwide, and in Australia it is a major cause of death in children under five, who most often drown in home swimming pools. Drowning can occur in pools, beaches, rivers, bathtubs, and buckets of water, and it can happen in seconds and in very shallow water. Prevention through supervision, barriers, swimming skills, and safety measures is the most important strategy, and immediate first aid including rescue breathing and CPR saves lives.

\section*{Epidemiology}
Around 250--300 people drown in Australia each year, and drowning is one of the leading causes of accidental death in children. Children under five are at the highest risk, with home swimming pools the most common site, and adults most often drown in rivers, beaches, and rock fishing. Males drown at around twice the rate of females. Non-fatal drowning can cause brain injury and long-term disability. Drowning rates have fallen with pool fencing laws, swimming education, and public education, but remain preventable.

\section*{Aetiology and Pathophysiology}
\begin{itemize}
    \item \textbf{Mechanism:} water enters the airway, blocking breathing and causing oxygen deprivation
    \item \textbf{Laryngospasm:} the airway closes in response to water, causing asphyxia
    \item \textbf{Hypoxia:} lack of oxygen damages the brain and organs, which determines the outcome
    \item \textbf{Supervision gaps:} most child drownings occur when supervision lapses for a short time
    \item \textbf{Risk factors:} lack of barriers around pools, no adult supervision, poor swimming skills, alcohol, and unsafe conditions
    \item \textbf{Rescue:} prompt removal from water and early CPR are critical to survival
\end{itemize}

\section*{Clinical Features}
\begin{itemize}
    \item A person found in or under water, unresponsive or struggling
    \item Coughing, difficulty breathing, or gasping
    \item Cold, blue skin and lips
    \item Drowsiness, confusion, or unconsciousness
    \item Vomiting and fluid in the lungs
    \item Cardiac arrest in severe cases
    \item Delayed effects: breathing problems and brain injury can develop over hours
\end{itemize}

\section*{Differential Diagnosis}
\begin{itemize}
    \item Drowning versus near-drowning, where the person survives the initial event
    \item Other causes of collapse in the water, including cardiac arrest, seizures, and hypothermia
    \item Other causes of breathing difficulty after water exposure, such as water inhalation without submersion
\end{itemize}

\section*{When to Seek Medical Care}
Anyone who has been rescued from the water should receive first aid immediately and be assessed in hospital, even if they seem fine, because complications such as breathing problems and brain swelling can develop later.

\textbf{Call 000 (or local emergency services) for:}
\begin{itemize}
    \item Any drowning incident, even if the person appears recovered
    \item A person who is unresponsive, not breathing, or gasping
    \item Difficulty breathing, coughing up fluid, or drowsiness after water exposure
    \item Hypothermia or signs of serious injury
\end{itemize}

\section*{Diagnosis and Investigations}
\begin{itemize}
    \item \textbf{Assessment of airway, breathing, and circulation:} the first priority in the emergency setting
    \item \textbf{Oxygen monitoring:} pulse oximetry and blood gases
    \item \textbf{Chest X-ray:} detects fluid in the lungs
    \item \textbf{Neurological assessment:} brain function and consciousness
    \item \textbf{Blood tests:} organ function and electrolytes
    \item \textbf{Hospital observation:} monitoring for delayed complications
\end{itemize}

\section*{Management}
\begin{itemize}
    \item \textbf{Rescue:} remove the person from the water safely
    \item \textbf{CPR:} start rescue breathing and chest compressions; for children, breathing support is especially important
    \item \textbf{Call for help:} call 000 and continue CPR until help arrives
    \item \textbf{Hospital care:} oxygen, breathing support, and monitoring
    \item \textbf{Treatment of complications:} lung injury, brain swelling, hypothermia, and infection
    \item \textbf{Supportive care:} warming, fluids, and intensive care as needed
    \item \textbf{Follow-up:} assessment and rehabilitation for survivors
\end{itemize}

\section*{Complications}
\begin{itemize}
    \item Death
    \item Brain injury from oxygen deprivation, causing permanent disability
    \item Lung damage and pneumonia
    \item Hypothermia
    \item Psychological trauma for the survivor, family, and rescuers
    \item Long-term rehabilitation needs
\end{itemize}

\section*{Prognosis and Outcomes}
The outcome of drowning depends mainly on how quickly the person is rescued and how soon breathing and circulation are restored. People who receive early CPR and regain breathing quickly usually recover fully, while prolonged oxygen deprivation can cause brain injury or death. Every minute of hypoxia worsens the outlook, so immediate action saves lives. Prevention is the most effective strategy, and most drownings are preventable.

\section*{Prevention and Self-Care}
\begin{itemize}
    \item Supervise children constantly near water; supervision should be close, continuous, and undistracted
    \item Fence home swimming pools with compliant fencing and self-closing gates
    \item Teach children swimming and water safety skills
    \item Learn CPR; refresher courses are widely available
    \item Wear life jackets for boating and rock fishing
    \item Never swim alone, and supervise others
    \item Avoid alcohol around water
    \item Empty buckets, wading pools, and bathtubs after use
\end{itemize}

\section*{Sources and Further Reading}
The following reputable sources provide further information for healthcare students and patients seeking reliable, current advice about drowning prevention.

\begin{itemize}
    \item Royal Life Saving Australia. \textit{Drowning prevention and water safety.} https://www.royallifesaving.com.au/
    \item Australian Resuscitation Council. \textit{CPR and drowning rescue guidelines.} https://resus.org.au/
    \item Kidsafe Australia. \textit{Home pool safety.} https://www.kidsafe.com.au/
    \item Surf Life Saving Australia. \textit{Beach safety information.} https://sls.com.au/
    \item World Health Organization. \textit{Drowning fact sheet.} https://www.who.int/
    \item Healthdirect Australia. \textit{Drowning and first aid.} https://www.healthdirect.gov.au/
\end{itemize}
